\section{Discussion}
\label{s:discussion}

\paragraph{LLMs interpret intent well but should not decide license facts alone.}
RQ1 shows a clear pattern: the recommender handles broad requests almost perfectly but degrades as requests become specific, exactly where developers most need help. The errors are not random: a few requirements (trademark protection, change documentation, copyleft level) account for most violations, and the model gravitates towards popular licenses such as Apache-2.0. Because these are factual properties already encoded in the license labels, a deterministic step can check them. This supports the neuro-symbolic design planned for Phase~2: the LLM converts the description into structured requirements, and a rule engine over the labels (and the OSADL compatibility matrix for dependencies) selects the licenses. The LLM then does what it does well, interpreting language, and the lookup, which cannot misstate a license's labels, does the matching. An ``oracle-input'' experiment, giving Stage~1 the exact requirements instead of the description, would separate misunderstanding errors from license-knowledge errors and indicate how much of the gap such a design can close.

\paragraph{An off-the-shelf verifier is cautious, which suits this domain.}
RQ2 shows that DeBERTa rarely confirms a claim that is not supported (precision 0.94), and it never confirmed a false claim in the human-labelled set. In licensing, the cost of a wrong ``verified'' is higher than the cost of a missing one, so this error profile is acceptable for a first version, provided that the interface presents ``unverified'' as ``not confirmed by the clause'' rather than ``wrong''. The verifier's weakness is concentrated in claims that require a small inferential step beyond the clause wording (``sell'' $\rightarrow$ commercial use; ``license the entire work under this License'' $\rightarrow$ same-license copyleft). Two remedies follow: prompting the explainer to phrase claims in the license's own vocabulary, which our clause-cited design already encourages, and fine-tuning the verifier on license-domain pairs such as ours, split by license to avoid leakage.

\paragraph{Grounding makes explanations checkable.}
Requiring each claim to cite and quote a specific clause made every explanation in our runs traceable: all quotes were verbatim, and each claim can be inspected by the user in one click. Even when the verifier is uncertain, the user can read the evidence directly, which is the transparency property we prioritised.

\paragraph{Expert-free evaluation is feasible with care.}
Computing answer keys from curated labels, generating test inputs with one model and validating them with another, and checking automatic labels with a small human sample gave us evaluation data whose quality we could measure (97.9\% label accuracy; $\kappa = 0.88$). The approach transfers to other domains where curated structured data exists but expert annotators do not.
